\section{Arc coordinates and periodic-orbit matching}\label{sec:arcs}

We first study a general profile
\(\phi\in C^\infty((0,\infty))\) and the scalar equation
\begin{equation}\label{eq:arc-ode}
  \frac{du}{dy}=\phi(u)-y.
\end{equation}
Here \(y\) is an endpoint coordinate, rather than physical time.
For each \(t>0\), let \(u(y;t)\) be the solution through
\[
  u(\phi(t);t)=t.
\]
The parameter \(t\) will be the peak height. At a lower height \(h\), the
two endpoint values of \(y\) determine a half-width \(r\) and a midpoint
\(M\), measured relative to \(\phi(h)\).

Figure~\ref{fig:arc-coordinates} illustrates these coordinates.
Our first task is to show that width parametrizes every positive-height
arc. We then continue the transverse arcs to height zero, where
matching the two endpoint pairs counts the periodic orbits.

\subsection{Global endpoint data and their transport}

\begin{lemma}[Global arc coordinates]\label{lem:arc-data}
For every \(0<h<t\), the solution through
\((y,u)=(\phi(t),t)\) has exactly two intersections with \(u=h\).
They are transverse and satisfy
\[
  y_<(h,t)<\phi(h)<y_>(h,t).
\]
The part of the solution between them has a unique maximum \(u=t\).
Set
\begin{equation}\label{eq:arc-endpoint-data}
  r(h,t)=\frac{y_>(h,t)-y_<(h,t)}2,\qquad
  M=\frac{y_>(h,t)+y_<(h,t)}2-\phi(h).
\end{equation}
At fixed \(h\), the lower endpoint decreases strictly with \(t\), the
upper endpoint increases strictly, and both derivatives are nonzero.
The map \(t\mapsto r(h,t)\) is a smooth increasing bijection from
\((h,\infty)\) onto \((0,\infty)\).

Consequently \(M\) is a smooth function of \((h,r)\in(0,\infty)^2\).
Writing \(v=M_r\), one has
\begin{equation}\label{eq:arc-strict-data}
  |M(h,r)|<r,\qquad |v(h,r)|<1,\qquad
  r^2\ge 2(t-h).
\end{equation}
For a fixed width, the corresponding base height is strictly increasing
with peak height: \(\partial h/\partial t>0\).

If \(\phi\) is smooth on the whole real line, the same conclusions hold
for arbitrary real \(h<t\). In that case, at each fixed peak \(t\), letting
\(h\) decrease from \(t\) to \(-\infty\) parametrizes the full range
\(0<r<\infty\).
\end{lemma}

\begin{proof}
Differentiating Equation~\eqref{eq:arc-ode} gives
\[
  u_{yy}=\phi'(u)u_y-1.
\]
Thus every critical point in the profile domain is a nondegenerate
maximum, with \(u_{yy}=-1\). On either side of the prescribed peak the
height therefore decreases strictly as one moves away from it, until
the level \(h\) is reached. A second critical point on either branch
would require a minimum between two maxima, which is impossible.

These intersections occur at finite \(y\). Indeed, while a branch remains
in \([h,t]\), the profile is bounded there. For sufficiently large
positive \(y\), Equation~\eqref{eq:arc-ode} gives \(u_y<-1\); for
sufficiently large negative \(y\), it gives \(u_y>1\). In the respective
outward directions these inequalities force the branch to reach \(h\).
There is also no finite-\(y\) failure of continuation while \(u\) stays
in the compact interval \([h,t]\). The strict derivative signs at the
intersections give the displayed endpoint inequalities.

All the endpoints depend smoothly on \((h,t)\), by smooth dependence
of the differential equation and transversality; see, for example,
\cite[Sections~2.2--2.3]{Perko} for the standard differential-equation
facts. To find their peak
derivatives, let \(\xi(y;t)=\partial_t u(y;t)\), with \(y\) held fixed.
Differentiating the moving initial condition gives
\(\xi(\phi(t);t)=1\), because \(u_y=0\) at the peak. Moreover,
\[
  \xi_y=\phi'(u)\xi,
\]
so \(\xi>0\) throughout the arc. Differentiating the endpoint equations
now gives
\[
  (y_<)_t=-\frac{\xi(y_<;t)}{r-M}<0,\qquad
  (y_>)_t=\frac{\xi(y_>;t)}{r+M}>0.
\]
In particular \(r_t>0\). After changing from \(t\) to \(r\),
\[
  v=\frac{(y_>)_t+(y_<)_t}{(y_>)_t-(y_<)_t}\in(-1,1).
\]
The other strict inequality, \(|M|<r\), follows directly from the
endpoint inequalities.

At fixed \(t\), differentiation with respect to \(h\) gives
\[
  (y_<)_h=\frac1{r-M},\qquad
  (y_>)_h=-\frac1{r+M}.
\]
Hence
\begin{equation}\label{eq:arc-width-height}
  \left.\frac{dr}{dh}\right|_t=-\frac{r}{r^2-M^2},\qquad
  \left.\frac{dM}{dh}\right|_t
       =\frac{M}{r^2-M^2}-\phi'(h).
\end{equation}
The first identity implies \(d(r^2)/dh\le-2\).
The nondegenerate maximum gives \(r\to0\) as \(t\downarrow h\), or as
\(h\uparrow t\) with the peak fixed. Integrating the inequality yields
\(r^2\ge2(t-h)\). At fixed \(h\), it follows that \(r\to\infty\) as
\(t\to\infty\). This proves the asserted bijection and its smooth
inverse. Also, at fixed width,
\[
  \left.\frac{\partial h}{\partial t}\right|_r
    =-\frac{r_t}{r_h}>0.
\]
The same proof applies to any compact interval \([h,t]\) in the domain
of a smooth profile. For a profile on the whole real line, the bound
\(r^2\ge2(t-h)\) also proves the last assertion as \(h\to-\infty\).
\end{proof}

\begin{figure}[tb]
\centering
\begin{tikzpicture}[x=1.4cm,y=1.05cm,>=Stealth,font=\small]
  \draw[->] (-2.4,0)--(3.05,0) node[right] {$y$};
  \draw[->] (-2.25,-.1)--(-2.25,3.55) node[above] {$u$};
  \draw[densely dashed,gray] (-2.25,.85)--(2.65,.85);
  \node[left] at (-2.25,.85) {$h$};
  \draw[thick] (-1.65,.85) .. controls (-1.15,2.50) and (-.45,3.0) .. (.30,3.0)
               .. controls (1.08,3.0) and (1.85,2.45) .. (2.45,.85);
  \fill (-1.65,.85) circle (1.5pt);
  \fill (2.45,.85) circle (1.5pt);
  \fill (.30,3.0) circle (1.5pt);
  \draw[densely dotted] (.30,3.0)--(-2.25,3.0);
  \node[left] at (-2.25,3.0) {$t$};
  \node[above] at (.30,3.0) {$(\phi(t),t)$};
  \draw[densely dotted] (-1.65,0)--(-1.65,.85);
  \draw[densely dotted] (2.45,0)--(2.45,.85);
  \node[below] at (-1.65,0) {$y_<$};
  \node[below] at (2.45,0) {$y_>$};
  \draw[densely dotted] (.4,.85)--(.4,-.55);
  \draw[<->] (-1.65,.53)--(.4,.53) node[midway,below] {$r$};
  \draw[<->] (.4,.53)--(2.45,.53) node[midway,below] {$r$};
  \draw[densely dotted] (-.45,.85)--(-.45,-.55);
  \node[below] at (-.45,-.65) {$\phi(h)$};
  \node[below] at (.95,-.65) {$\phi(h)+M$};
  \draw[thin] (.4,-.55)--(.95,-.65);
  \draw[<->] (-.45,-.42)--(.4,-.42) node[midway,above] {$M$};
\end{tikzpicture}
\caption{An arc at base height $h$ and peak height $t$, drawn
schematically with $M>0$. Its endpoint half-width is $r$; its midpoint
is measured relative to $\phi(h)$. Both endpoints move outwards as
the peak rises, giving $|M_r|<1$. The two physical half-orbits form
a periodic orbit precisely when their height-zero endpoint pairs
agree, as proved in Proposition~\ref{prop:cycle-matching}.}
\label{fig:arc-coordinates}
\end{figure}


Put
\[
  k(h)=\phi'(h),\qquad q=r^2-M^2,\qquad
  \alpha=\frac rq,\qquad
  D=\partial_h-\alpha\partial_r.
\]
By Equation~\eqref{eq:arc-width-height}, \(D\) differentiates along a
fixed-peak arc as its base height increases. We call these curves in
the \((h,r)\)-plane \emph{characteristics}. Their basic transport
identities are
\begin{equation}\label{eq:arc-transport}
  DM=\frac Mq-k(h),\qquad
  Dv=-\frac{2Mr(1-v^2)}{q^2}.
\end{equation}
For clarity, the second identity includes a commutator term. For any
smooth function \(f(h,r)\),
\[
  D(f_r)=(Df)_r+\alpha_r f_r,\qquad
  \alpha_r=\frac1q-\frac{2r(r-Mv)}{q^2}.
\]
Using the first identity in Equation~\eqref{eq:arc-transport} therefore
gives
\[
  Dv=\frac{2v}{q}
       -\frac{2(M+rv)(r-Mv)}{q^2}
     =-\frac{2Mr(1-v^2)}{q^2},
\]
as claimed. All partial \(r\)-derivatives here are taken at fixed base
height.

Introduce the angles
\[
  \theta=\operatorname{arctanh}(M/r),\qquad
  L=\operatorname{arctanh}(v).
\]
When a characteristic is parametrized by increasing width \(r\), a
prime denotes the total derivative along it; thus \(f'=-Df/\alpha\).
Equation~\eqref{eq:arc-transport} first gives
\[
  M'=-\frac Mr+\frac qr k(h).
\]
Substituting \(M=r\tanh\theta\) and \(v=\tanh L\) then yields
\begin{equation}\label{eq:arc-characteristics}
\begin{aligned}
  h'&=-r\operatorname{sech}^2\theta,\\
  \theta'&=k(h)-\frac{\sinh(2\theta)}r,&
  L'&=\frac{\sinh(2\theta)}r,\\
  k(h)'&=-r\operatorname{sech}^2\theta\,\phi''(h).
\end{aligned}
\end{equation}
In particular,
\begin{equation}\label{eq:arc-height-bound}
  t-\frac{r^2}{2}\le h\le t.
\end{equation}
The initial values of the angles at width zero, and the corresponding
integral formula for \(L\), will follow from the local scaling below.

\subsection{Profile variation and smooth small-width data}

\begin{lemma}[Variation at fixed height and width]\label{lem:variation}
Consider a smooth one-parameter family of profiles \(\phi_\eta\), with
\(\phi_0=\phi\), and fix a regular base height \(h\) and a width \(r>0\).
All profiles need only be defined on a neighborhood of the relevant
compact arcs. Write
\[
  \psi(u)=\left.\partial_\eta\phi_\eta(u)\right|_{\eta=0},
  \qquad Y=M+\phi(h),\qquad
  K(y)=\exp\left(-\int_{y_<}^{y}\phi'(u(z))\,dz\right).
\]
Then the first variation \(\delta Y\), with \(h,r\) fixed, satisfies
\begin{equation}\label{eq:arc-variation}
 \bigl[(M+r)K(y_>)-(M-r)K(y_<)\bigr]\delta Y
       =\int_{y_<}^{y_>}K(y)\psi(u(y))\,dy.
\end{equation}
The coefficient of \(\delta Y\) is strictly positive. The variation of
the relative midpoint is \(\delta M=\delta Y-\psi(h)\).

Consequently, if two profiles smooth on \((0,\infty)\) satisfy, at
a fixed \(h>0\),
\[
  \phi_1(u)-\phi_1(h)>\phi_0(u)-\phi_0(h)
  \qquad\text{for every }u>h,
\]
then \(M_1(h,r)>M_0(h,r)\) for every \(r>0\). The analogous
comparison holds at any real base height for two entire smooth
profiles.
\end{lemma}

\begin{proof}
Both endpoint coordinates have variation \(\delta Y\), since their
difference \(2r\) is fixed. Let \(\delta u\) denote variation at fixed
\(y\). Differentiating the differential equation and endpoint
conditions gives
\[
  (\delta u)_y=\phi'(u)\delta u+\psi(u),\qquad
  \delta u(y_<)=(M-r)\delta Y,\quad
  \delta u(y_>)=(M+r)\delta Y.
\]
Multiplication by \(K\) and integration prove
Equation~\eqref{eq:arc-variation}. Its coefficient is positive because
\(K>0\), \(M+r>0\), and \(M-r<0\).

For the comparison, first translate each endpoint coordinate by its
own base value \(\phi_i(h)\). This replaces its profile by
\(\phi_i(u)-\phi_i(h)\), without changing \(M_i\). Interpolate linearly
between the translated profiles. The variation vanishes at the base
and is strictly positive in every arc interior. Equation~\eqref{eq:arc-variation}
then makes the derivative of the midpoint strictly positive along
the interpolation. Integrating in the interpolation parameter proves
the comparison.
\end{proof}

We will use three exact variations at the zero profile. The entire
profile version of Lemma~\ref{lem:arc-data} permits base height zero.
At width \(w>0\), the arc is
\[
  u_w(y)=\frac{w^2-y^2}{2},\qquad -w\le y\le w,
\]
with \(M=0\) and \(K=1\). Thus a profile variation \(\psi\) with
\(\psi(0)=0\) gives
\[
  \delta M(0,w)=\frac1{2w}\int_{-w}^{w}\psi(u_w(y))\,dy.
\]
Direct polynomial integration gives
\[
 \int_{-w}^{w}u_w\,dy=\frac{2w^3}{3},\qquad
 \int_{-w}^{w}\frac{u_w^2}{2}\,dy=\frac{2w^5}{15},\qquad
 \int_{-w}^{w}\frac{u_w^3}{6}\,dy=\frac{2w^7}{105}.
\]
Smooth dependence allows differentiation of these first variations
with respect to the positive width \(w\). Therefore
\begin{equation}\label{eq:parabolic-variations}
\begin{array}{c|cc}
  \psi(s)&\delta M(0,w)&\delta M_w(0,w)\\ \hline
  s&w^2/3&2w/3\\
  s^2/2&w^4/15&4w^3/15\\
  s^3/6&w^6/105&6w^5/105
\end{array}
\end{equation}
These constants will determine the local quadratic fit.

\begin{lemma}[Smooth scaling at zero width]\label{lem:arc-scaling}
Near any fixed positive base height, the small-width arc data have
smooth parameter-dependent expansions
\begin{equation}\label{eq:arc-small-width}
\begin{aligned}
  M(h,r)&=\frac13 k(h)r^2+r^3 A(h,r),\\
  v(h,r)&=\frac23 k(h)r+r^2 B(h,r),\\
  t(h,r)&=h+r^2\left(\frac12+r\,C(h,r)\right),
\end{aligned}
\end{equation}
where the remainder functions extend smoothly through \(r=0\).
The same assertion holds jointly for any finite-dimensional smooth
family of profiles. In particular, these expansions may be
differentiated, and \(\left.\partial h/\partial t\right|_r\to1\)
as \(r\downarrow0\).
The assertions also hold near any regular base height of a profile
defined on an open interval containing that height.
\end{lemma}

\begin{proof}
For positive \(r\), introduce the scaled height and endpoint coordinate
\[
  s=\frac{u-h}{r^2},\qquad
  \zeta=\frac{y-\phi(h)}r.
\]
The scaled equation is
\[
  \frac{ds}{d\zeta}=g(s)-\zeta,\qquad
  g(s)=\frac{\phi(h+r^2s)-\phi(h)}r.
\]
Its relative midpoint at scaled width \(1\) is \(M(h,r)/r\).
More generally, if the scaled profile is held fixed while its width
\(w\) varies near \(1\), its midpoint is \(M(h,rw)/r\). Its width
derivative at \(w=1\) is therefore \(v(h,r)\).

To justify smoothness through zero, introduce a signed auxiliary
parameter \(\rho\), while keeping the scaled width positive and near
\(1\). Define
\begin{equation}\label{eq:arc-scaling-profile}
  g(h,\rho,s)
   =\rho s\int_0^1\phi'(h+\tau\rho^2s)\,d\tau.
\end{equation}
Choose a fixed compact scaled-height interval containing a neighborhood
of \([0,1/2]\), and restrict \(h\) to a neighborhood of the given
regular base height. For sufficiently small \(|\rho|\), all arguments
of \(\phi'\) lie in its smooth domain. Equation~\eqref{eq:arc-scaling-profile}
is then a smooth family through \(\rho=0\), equals the preceding
difference quotient when \(\rho\ne0\), and is the zero profile at
\(\rho=0\). Negative \(\rho\) is only a parameter in this family;
no negative physical width is used.

At the zero profile, an arc of scaled peak height \(T\) has endpoints
\(\pm\sqrt{2T}\), and its width derivative with respect to \(T\) equals
\(1\) at \(T=1/2\). Smooth dependence of the scaled differential
equation, endpoint transversality, and the implicit function theorem
therefore give smooth peak, midpoint, and midpoint-width-derivative
data at scaled width \(1\) for the family in
Equation~\eqref{eq:arc-scaling-profile}. All of this uses only a
neighborhood of the compact parabolic arc, so the local profile domain
is sufficient.

The derivative of \(g\) with respect to \(\rho\) at zero is \(k(h)s\).
Equation~\eqref{eq:parabolic-variations} consequently gives first
derivatives \(k(h)/3\) and \(2k(h)/3\) for the scaled midpoint and its
width derivative. Their values at \(\rho=0\) are zero. The scaled peak
is \(1/2+O(\rho)\). Taylor's formula with a smooth remainder now gives
Equation~\eqref{eq:arc-small-width} on setting \(\rho=r>0\) and
undoing the scaling. It also proves the stated smoothness of the
remainder functions, including any smooth profile parameters.
Finally \(t_h=1+O(r^2)\) at fixed \(r\), so the inverse derivative
\(\partial h/\partial t\) tends to \(1\).
\end{proof}

On a characteristic with fixed smooth peak \(t\),
Lemma~\ref{lem:arc-scaling} and Equation~\eqref{eq:arc-height-bound}
give
\[
  \theta(r)=\frac13 k(t)r+O(r^2),\qquad
  L(r)=\frac23 k(t)r+O(r^2).
\]
Both angles therefore tend to zero, and their transport equation
integrates to
\begin{equation}\label{eq:arc-L-integral}
  L(r)=\int_0^r\frac{\sinh(2\theta(s))}{s}\,ds.
\end{equation}
The integrand is bounded near zero. All these expansions are smooth
in finite profile parameters. In particular, for a quadratic profile
whose slope and curvature at the peak are \(z\) and \(\kappa\), let
\(j\) be its slope at the moving base. Translating the peak to height
zero and applying the local form of Lemma~\ref{lem:arc-scaling} yields,
locally uniformly for finite \((z,\kappa)\),
\begin{equation}\label{eq:arc-peak-expansions}
\begin{aligned}
  \theta(s,z,\kappa)&=\frac13 zs+O(s^2),&
  \theta_z(s,z,\kappa)&=\frac13 s+O(s^2),\\
  L(s,z,\kappa)&=\frac23 zs+O(s^2),&
  j(s,z,\kappa)&=z+O(s^2),\qquad j_z\longrightarrow1 .
\end{aligned}
\end{equation}
Here the differentiated statements follow from smooth remainders,
rather than from differentiating an unspecified error bound.

\subsection{Continuation to a nonsmooth boundary height}

For the next lemma, assume additionally that \(\phi\) extends
continuously to \([0,\infty)\) and that \(\phi(0)=0\).
Derivatives of \(\phi\) need not extend to zero.

\begin{lemma}[Transverse extension to the axis]\label{lem:axis-extension}
Each peak \(t>0\) has finite limiting endpoints
\[
  a(t)=\lim_{h\downarrow0}y_<(h,t)\le0,\qquad
  b(t)=\lim_{h\downarrow0}y_>(h,t)\ge0.
\]
The set
\[
  \mathcal T_\phi=\{t>0:a(t)<0<b(t)\}
\]
is either empty or an open ray in \((0,\infty)\). If it is nonempty,
\(a,b\) are smooth there, with \(a_t<0<b_t\), and
\[
  r(0,t)=\frac{b(t)-a(t)}2
\]
is a smooth increasing bijection onto an open width ray
\(\mathcal I_\phi\subset(0,\infty)\). If \(\mathcal T_\phi\) is empty,
set \(\mathcal I_\phi=\varnothing\).

On \(\mathcal I_\phi\), the limiting midpoint \(M(0,r)\) is smooth,
and
\[
  |M(0,r)|<r,\qquad |M_r(0,r)|<1.
\]
As \(h\downarrow0\), all partial \(r\)-derivatives of \(M(h,r)\)
converge locally uniformly on \(\mathcal I_\phi\) to the corresponding
derivatives of \(M(0,r)\).
\end{lemma}

\begin{proof}
At fixed peak, the endpoint derivatives in
Equation~\eqref{eq:arc-width-height} show that the lower endpoint
decreases and the upper endpoint increases as \(h\downarrow0\).
They remain bounded. To see this explicitly, put
\(B=\sup_{0\le u\le t}|\phi(u)|\). If an upper endpoint lies beyond
\(B\), integrate \(u_y\le B-y\) from \(B\) to that endpoint, using
\(u(B)\le t\). This gives \(y_>\le B+\sqrt{2t}\).
The inequality \(u_y\ge-B-y\) similarly gives
\(y_<\ge-B-\sqrt{2t}\). The remaining bounds follow from
\(y_<<\phi(h)<y_>\) and \(|\phi(h)|\le B\).
The finite limits thus exist, and the signs follow by continuity of
\(\phi\) at zero.

For two peaks \(t_2>t_1\), the positive-height endpoint order passes
to the weak limiting order
\[
  a(t_2)\le a(t_1),\qquad b(t_2)\ge b(t_1).
\]
Consequently \(\mathcal T_\phi\) is upward closed. We next show both
its openness and the local smoothness asserted in the lemma.

Fix \(t_0\in\mathcal T_\phi\). At a sufficiently small fixed positive
level \(\delta<t_0\), the two endpoint branches can be continued
downwards by
\begin{equation}\label{eq:arc-inverse-branch}
  \frac{dy}{du}=\frac1{\phi(u)-y}=:g(u,y).
\end{equation}
Choose compact neighborhoods of the limiting endpoints \(a(t_0)\)
and \(b(t_0)\) and make \(\delta\) smaller if necessary. On each
resulting branch rectangle, \(|\phi(u)-y|\) is bounded away from zero
for \(0\le u\le\delta\). The function \(g\) is continuous in \(u\),
smooth in \(y\), and has bounded derivatives of every fixed order
with respect to \(y\).

At \(u=\delta\), the branch initial values are smooth functions of
the peak \(t\) near \(t_0\). The integral equation for
Equation~\eqref{eq:arc-inverse-branch} and its parameter-variation
equations show that the continued branches and all their \(t\)
derivatives are continuous down to \(u=0\), locally uniformly in \(t\).
For the first derivative the formula is explicit:
\[
  y_t(u,t)=y_t(\delta,t)
      \exp\left(\int_\delta^u g_y(v,y(v,t))\,dv\right).
\]
The exponential factor is positive and bounded away from zero.
Thus the two nonzero variation signs persist uniformly near \(t_0\).
Higher derivatives satisfy linear variation equations with
continuous bounded coefficients and forcing built from lower
derivatives, which proves the same assertion inductively.
In particular, the strict endpoint signs persist for nearby peaks.
This proves openness. An open upward-closed subset of \((0,\infty)\)
is an open ray, or is empty.

On that ray, \(r_t=(b_t-a_t)/2>0\). Passing to the limit in
\(r^2\ge2(t-h)\) gives \(r(0,t)^2\ge2t\); hence the width tends to
infinity as the peak does. Its image is therefore one open ray,
and its inverse is smooth.

It remains to justify convergence at fixed width, rather than only
at fixed peak. Near \(t_0\), the preceding bounds give
\(r_t(h,t)\ge c_0>0\) uniformly for all sufficiently small
\(h\ge0\). Uniform strict monotonicity brackets the inverse peak
\(t=T(h,r)\) in a common compact interval for nearby widths.
This inverse is continuous in \((h,r)\), smooth in \(r\), and
\[
  T_r=\frac1{r_t(h,T(h,r))}.
\]
Higher \(r\)-derivatives are rational expressions in the continuous
peak derivatives, with powers of \(r_t\) in the denominator.
They consequently converge locally uniformly as \(h\downarrow0\).
Composing the endpoint data with \(T(h,r)\), and subtracting
\(\phi(h)\) from their midpoint, proves the claimed convergence of
all \(r\)-derivatives. No derivative of \(\phi\) at zero is needed.
Finally the strict endpoint signs give \(|M|<r\), and the retained
opposite peak derivatives give \(|M_r|<1\), exactly as in
Lemma~\ref{lem:arc-data}.
\end{proof}

\subsection{Matching half-arcs and isolated periodic orbits}

We now connect the general construction to a smooth classical
Liénard system. This step requires only \(F\in C^\infty(\mathbb R)\)
and \(F(0)=0\), not a degree restriction. Define
\[
  \phi_\pm(u)=F(\pm\sqrt{2u}),\qquad u\ge0,
\]
and let \(M_\pm\) and \(\mathcal I_\pm\) be their midpoint data and
transverse width domains from Lemma~\ref{lem:axis-extension}.
These profiles are smooth for \(u>0\), continuous at zero, and vanish
there.

\begin{proposition}[Periodic-orbit matching]\label{prop:cycle-matching}
For the system
\begin{equation}\label{eq:arc-lienard-system}
  \dot x=y-F(x),\qquad \dot y=-x,
\end{equation}
the images of nonconstant periodic solutions are in one-to-one
correspondence with the zeros of
\begin{equation}\label{eq:cycle-matching}
  \Delta(r)=M_+(0,r)-M_-(0,r)
\end{equation}
on the single open interval
\(\mathcal I=\mathcal I_+\cap\mathcal I_-\).
If this interval is empty, there are no such periodic solutions.
Each periodic orbit has exactly one positive and one negative
intersection with the \(y\)-axis. Under the correspondence,
isolated periodic orbits are exactly isolated zeros of \(\Delta\).
\end{proposition}

\begin{proof}
On either open half-plane, \(y\) is strictly monotone in physical
time. With \(u=x^2/2\), division of the two time derivatives gives
\[
  \frac{du}{dy}=F(x)-y=\phi_\pm(u)-y.
\]
Thus any excursion in one half-plane is an arc of
Equation~\eqref{eq:arc-ode}.

Every nonconstant periodic solution has both signs of \(x\), since
integrating \(\dot y=-x\) over a period gives \(\int x\,dt=0\).
The only equilibrium is \((0,0)\), and uniqueness prevents a
nonconstant orbit from containing it. At any axis intersection,
\(\dot x=y\ne0\); the intersection is transverse. There are finitely
many in a period. Otherwise an accumulation time would give an
axis point at which either the crossing derivative is nonzero,
contradicting accumulation, or the solution reaches the equilibrium.

A positive axis intersection enters \(x>0\), and a negative one
enters \(x<0\). Between consecutive intersections the transformed
height is positive, vanishes at the endpoints, and has a unique
positive peak by Lemma~\ref{lem:arc-data}. Both endpoint signs are
strict, so the arc belongs to the transverse domain of
Lemma~\ref{lem:axis-extension}.

For either profile, increasing the positive endpoint strictly
decreases the negative endpoint, by peak order. The physical
right-half-plane map from positive to negative endpoints is
therefore strictly decreasing, and so is the physical left-half-plane
map from negative to positive endpoints. Their composition on
successive positive intersections is strictly increasing.
An increasing map cannot have a finite periodic sequence of
distinct points: if one iterate is larger, or smaller, than its
predecessor, strict order persists under further iteration.
Hence a closed orbit has one distinct positive intersection and
one distinct negative intersection, and consists of one arc of
each sign.

The two arcs of such an orbit have the same endpoints, so their
widths and midpoints agree. This gives a zero of
Equation~\eqref{eq:cycle-matching}. Conversely, a zero at width \(r\)
gives the same endpoints
\[
  a=M_\pm(0,r)-r<0,\qquad
  b=M_\pm(0,r)+r>0
\]
for the two arcs. Recover \(x\) as \(+\sqrt{2u_+(y)}\) on the right
and \(-\sqrt{2u_-(y)}\) on the left, and introduce physical time by
\(\dot y=-x\). On each open arc,
\(\dot x=-u_y=y-F(x)\), as required. The one-sided endpoint derivatives
are
\[
  (u_\pm)_y(a)=-a>0,\qquad (u_\pm)_y(b)=-b<0.
\]
Thus \(u_\pm\) has a simple zero at each endpoint and
\(dy/\sqrt{2u_\pm(y)}\) is integrable there. The physical travel
times are finite. The derivatives of the reconstructed solution
extend to the axis with values \((\dot x,\dot y)=(y,0)\).
The two arcs consequently join as a solution of the original smooth
vector field; uniqueness gives a closed orbit.

Different zeros yield different positive intersections, because
\[
  \frac{d}{dr}\bigl(r+M_+(0,r)\bigr)=1+M_{+,r}(0,r)>0.
\]
They therefore yield different geometric orbits. This proves the
bijection.

Finally fix a matched orbit, its positive intersection \((0,s_0)\),
and its first positive return time \(T>0\). On a short positive-axis
section \(S\) around \(s_0\), the first return ordinate \(P(s)\)
and time \(\tau(s)\) are smooth and satisfy \(\tau(s)\to T\) as
\(s\to s_0\). To see that these are the first returns, take disjoint
small time windows about the reference axis crossings. On the
intervening compact time segments, \(|x|\) is bounded away from
zero. Smooth flow dependence and transversality give exactly the
nearby crossings in those windows and no others.
The function \(s=r+M_+(0,r)\) is a local smooth coordinate with
positive derivative. In this coordinate,
\[
  \Delta(r)=0\quad\Longleftrightarrow\quad P(s)=s
\]
near the reference root.

If distinct roots tend to the reference root, the corresponding
section points tend to \(s_0\), their return times tend to \(T\),
and smooth flow dependence on a common compact time interval makes
their entire orbit images tend to the reference image. They are
distinct by the preceding monotonicity. Hence a nonisolated root
cannot give an isolated periodic orbit.

For the converse, suppose the root is isolated and shrink \(S\)
so that \(s_0\) is its only return fixed point. If \(\Phi_\tau\)
denotes the local flow, the map
\((s,\tau)\mapsto\Phi_\tau(0,s)\) is a local diffeomorphism near
each point \((s_0,\tau)\), \(0\le\tau\le T\), since the initial
section is transverse. Finitely many such flow neighborhoods cover
the compact reference orbit and give an open neighborhood \(U\)
whose every point lies on a trajectory meeting \(S\).
Any periodic orbit meeting \(U\) therefore meets \(S\). Its unique
positive-axis intersection must be a fixed point of the local
return map, and hence equals \(s_0\). Uniqueness then makes it the
reference orbit. This proves isolation, without an assumption on
hyperbolicity or a global bound on periods.
\end{proof}

\begin{remark}
Every periodic orbit in Proposition~\ref{prop:cycle-matching} lies
inside the open matching domain. A limiting endpoint at the
equilibrium is not an additional periodic orbit. Nonisolated zeros,
including a zero interval and its endpoints within the domain,
do not count as limit cycles. Isolated multiple zeros remain
covered by the same correspondence.
\end{remark}
